\documentclass[11pt,a4paper]{article}

\usepackage[utf8]{inputenc}
\usepackage[indonesian]{babel}
\usepackage[a4paper,top=2cm,bottom=2cm,left=2cm,right=2cm]{geometry}
\usepackage{amsmath,amsfonts,amssymb}
\usepackage{graphicx}
\usepackage{booktabs}
\usepackage{enumitem}
\usepackage{xcolor}
\usepackage{tcolorbox}
\usepackage{fancyhdr}
\usepackage{titlesec}
\usepackage{hyperref}

% Color Definitions
\definecolor{primaryblue}{RGB}{24, 43, 73}
\definecolor{secondaryblue}{RGB}{41, 128, 185}
\definecolor{accentblue}{RGB}{41, 128, 185}
\definecolor{boxbg}{RGB}{245, 248, 252}
\definecolor{goldaccent}{RGB}{211, 84, 0}

\pagestyle{fancy}
\fancyhf{}
\fancyhead[L]{\small\textbf{\color{primaryblue}Lembar Kerja Mahasiswa \& Problem Set (Minggu 14)}}
\fancyhead[R]{\small\color{gray}Analisis Deret Waktu -- Teknik Elektro}
\fancyfoot[L]{\footnotesize\color{gray}Program Studi Teknik Elektro FT UNIB}
\fancyfoot[R]{\footnotesize\thepage}

\titleformat{\section}{\color{primaryblue}\normalfont\large\bfseries}{\thesection.}{0.5em}{}[\titlerule]
\setlist{nosep}

\begin{document}

% HEADER
\begin{center}
    {\large \textbf{\color{primaryblue}UNIVERSITAS BENGKULU -- FAKULTAS TEKNIK}}\\[0.1cm]
    {\large \textbf{\color{secondaryblue}JURUSAN TEKNIK ELEKTRO}}\\[0.25cm]
    {\LARGE \textbf{\color{primaryblue}LEMBAR KERJA \& PROBLEM SET MINGGU 14}}\\[0.1cm]
    {\normalsize \textbf{Topik: Arsitektur Transformer dan Multi-Head Attention}}
\end{center}

\vspace{0.2cm}

\begin{tcolorbox}[colback=boxbg,colframe=primaryblue,arc=2mm,title=\textbf{Identitas Mahasiswa}]
\small
\begin{tabular}{p{2.5cm}p{5.5cm}p{2.5cm}p{4.5cm}}
\textbf{Nama} & : \dotfill & \textbf{NPM} & : \dotfill \\
\textbf{Kelas/Kelompok} & : \dotfill & \textbf{Tanggal} & : \dotfill \\
\end{tabular}
\end{tcolorbox}

\vspace{0.2cm}

\section{Panduan Pengerjaan}
Problem set ini dirancang berjenjang berdasarkan Taksonomi Bloom (Level C1 sampai C6) untuk mengukur capaian \textbf{Sub-CPMK 14}. Kerjakan secara mandiri, rapi, dan sistematis.

\section{Problem Set (Level C1 -- C6)}

\begin{enumerate}[leftmargin=1.8em, label=\textbf{Soal \arabic*.}]
    
    % C1 - Remembering
    \item \textbf{[C1 -- Mengingat / \textit{Remembering}]} \\
    Jelaskan secara konseptual tujuan dari mekanisme \textit{Self-Attention} pada model berbasis Arsitektur Transformer dalam konteks analisis deret waktu.
    \vspace{2.2cm}

    % C2 - Understanding
    \item \textbf{[C2 -- Memahami / \textit{Understanding}]} \\
    Apa fungsi modul \textit{Positional Encoding} dalam model Transformer? Mengapa Transformer membutuhkannya, tidak seperti RNN atau LSTM yang memproses data berurutan secara natural?
    \vspace{2.8cm}

    % C3 - Applying
    \item \textbf{[C3 -- Menerapkan / \textit{Applying}]} \\
    Dalam perhitungan attention $Attention(Q, K, V) = softmax(\frac{QK^T}{\sqrt{d_k}})V$, jika matriks \textit{Query} ($Q$) dan \textit{Key} ($K$) adalah input dari sequence historis daya yang sama sepanjang $T$, hitunglah secara konseptual ukuran dimensi (\textit{shape}) dari matriks dot product $QK^T$ (tanpa memperhatikan jumlah batch atau num\_heads).
    \vspace{2.8cm}

    \newpage

    % C4 - Analyzing
    \item \textbf{[C4 -- Menganalisis / \textit{Analyzing}]} \\
    Lakukan analisis terperinci keunggulan \textit{Multi-Head Attention} dibanding \textit{Single-Head Attention} untuk mendeteksi berbagai jenis pola periodisitas pada profil beban puncak secara bersamaan.
    \vspace{3.5cm}

    % C5 - Evaluating
    \item \textbf{[C5 -- Mengevaluasi / \textit{Evaluating}]} \\
    Evaluasilah tren peralihan para peneliti dari penggunaan algoritma LSTM ke model berbasis Transformer (misalnya PatchTST atau Informer) untuk memecahkan problem \textit{Long-Term Time Series Forecasting} (LTSF) di sistem kelistrikan pintar. Apakah Transformer selalu superior pada dataset skala menengah ke bawah?
    \vspace{3.5cm}

    % C6 - Creating
    \item \textbf{[C6 -- Merancang / \textit{Creating}]} \\
    Rancang struktur konseptual skema "Patching" pada deret waktu (seperti pada model PatchTST), di mana vektor univariat panjang dipecah-pecah menjadi banyak segmen (patch) yang \textit{overlapping} sebelum dimasukkan sebagai input ke layer encoder Transformer, dan jelaskan keuntungan skema tersebut untuk menangkap semantik struktur lokal data deret waktu listrik.
    \vspace{4.5cm}

\end{enumerate}

\end{document}
